\chapter{Kako se ovaj udžbenik koristi}
\label{ch:kako-koristiti}

\section{Čemu ovaj udžbenik služi}

Ovaj udžbenik ima jedan zadatak: da s minimalnim predznanjem možeš u cijelosti
razumjeti i obraniti seminarski rad iz kolegija \emph{Optimizacijske metode u razvoju i
upravljanju sustava}. Ne pretpostavlja da već znaš što je euklidski prostor, što znači
$\langle x, y \rangle$ ili zašto se piše $\coloneqq$ umjesto $=$. Sve se to definira
prije prve upotrebe, u Poglavlju~\ref{ch:rjecnik}.

Građen je \textbf{po zadacima}, ne po temama. Ideš zadatak po zadatak, podzadatak po
podzadatak, točno onim redom kojim ćeš ih morati izložiti na obrani. Teorija ulazi
onda kad zatreba i tamo gdje zatreba --- ne prije.

\section{Četiri boje: odakle koji tekst dolazi}

Najveći problem pri učenju iz vlastitih bilježaka jest da se poslije ne zna
\emph{tko je što rekao}: je li ovo profesorova definicija, je li ovo bio dio zadatka,
ili sam to sâm zaključio. Zato u ovom udžbeniku \textbf{boja označava isključivo
podrijetlo teksta}, ništa drugo.

\begin{izvor}{Predavanje 2 --- Definicije, klasifikacija, konveksnost, slajd 60/71 (PDF str.~61)}
Crveni okvir je doslovno ono što piše na profesorovu slajdu. Stranica slajda umeće se
kao slika, pa ispod nje slijedi pretipkani sadržaj s objašnjenjem svakog simbola.
Referenca uvijek ima \emph{dva} broja: broj slajda ispisan u podnožju samog slajda i
broj stranice u PDF-u. Ta se dva broja razlikuju, jer beamer svaki slajd s postupnim
otkrivanjem sprema kao više PDF stranica.

\slajd{p02-konveksna-funkcija}
\end{izvor}

\begin{zadatak}{Zadatak 1, uvodna formulacija}
Sivi okvir je doslovni tekst zadatka --- izravno sa skena koji je profesor potpisao.
Unutar sivog okvira \textbf{nema} ničega što je netko poslije dodao, zaključio ili
protumačio. Ako nešto u zadatku nije rečeno, u sivom okviru toga neće biti.
\skenzad[0.98]{z1-original}
\end{zadatak}

\begin{seminar}{Seminar v1.2, §1.1}
Plavi okvir je doslovni izvadak iz predanog seminarskog rada v1.2 --- ista notacija,
iste jednadžbe, isti brojevi kao u PDF-u koji je predan. Tako uvijek znaš što točno
braniš.
\end{seminar}

\begin{pojam}{Konveksna funkcija}
Žuti okvir je definicija ili formula koju treba znati napamet. Numeriran je, pa se na
njega poslije referira brojem.

Funkcija $f: S \to \R$ je \textbf{konveksna} ako je $S$ konveksan skup i ako za sve
$x_1, x_2 \in S$ te za svaki $\alpha$, $0 < \alpha < 1$, vrijedi
\begin{equation}
f\bigl(\alpha x_1 + (1-\alpha) x_2\bigr) \;\leq\; \alpha f(x_1) + (1-\alpha) f(x_2).
\end{equation}
\end{pojam}

\section{Tihi elementi}

Sve što nije jedan od ta četiri okvira nema ispunu ni boju --- samo tanku crtu s lijeve
strane. To je namjerno: da četiri boje ostanu jedini glasan signal na stranici.

\begin{predznanje}
\begin{itemize}[nosep,leftmargin=*]
  \item Popis simbola i pojmova koje moraš znati prije nego kreneš dalje.
  \item Uz svaki stoji stranica na kojoj je objašnjen.
\end{itemize}
\end{predznanje}

\begin{korak}
Numerirani izvod. Svaki korak imenuje simbole koje uvodi; nijedan algebarski korak se
ne preskače uz \emph{očito slijedi}.
\begin{enumerate}[nosep,leftmargin=*]
  \item Zapiši ono što znaš.
  \item Uvrsti.
  \item Zaključi, i reci \emph{zašto} smiješ zaključiti.
\end{enumerate}
\end{korak}

\begin{zamka}
Greška koja se stvarno radi, i objašnjenje zašto je pogrešna. Nije popis općenitih
savjeta nego konkretan promašaj koji te na obrani može stajati ocjene.
\end{zamka}

\begin{ploca}
Kako nacrtati sliku rukom, na ploči, bez računala: koje osi, kojim redom, koje točke
označiti i što napisati uz njih.
\end{ploca}

\begin{izgovor}
Doslovne rečenice koje možeš izgovoriti pred ispitivačem. Nisu parafraza --- napisane
su tako da ih se može reći od riječi do riječi.
\end{izgovor}

\begin{pitanje}{Zašto je Hesseova matrica ovdje konstantna?}
Okvir s pitanjem ispitivača, a odmah ispod modelski odgovor.

\textbf{Odgovor.} Jer je funkcija cilja kvadratna: druge parcijalne derivacije
kvadratne funkcije su konstante, pa Hesseova matrica ne ovisi o $x$.
\end{pitanje}

\section{Ritam svakog podzadatka}

Svaki podzadatak --- 1a, 1b, \dots, 4c --- ide istim redom od devet koraka. Kad se
navikneš na taj ritam, znaš gdje ćeš što naći bez traženja.

\begin{enumerate}[leftmargin=*]
  \item \textbf{Što se traži} --- jedna rečenica, bez žargona.
  \item \textbf{Sivi okvir} --- izvorni tekst zadatka.
  \item \textbf{Predznanje} --- što moraš znati prije, sa stranicama.
  \item \textbf{Crveni okvir} --- slajd s predavanja ili vježbi, pa transkripcija.
  \item \textbf{Žuti okviri} --- definicije i formule koje se koriste.
  \item \textbf{Izvod korak po korak} --- bez preskakanja.
  \item \textbf{Plavi okvir} --- kako je riješeno u seminaru.
  \item \textbf{Graf} --- slika, čitanje slike, crtanje rukom.
  \item \textbf{MATLAB i stvarni ispis} --- kod i doslovni ispis programa.
\end{enumerate}

Na kraju svakog zadatka slijedi modul za obranu: pitanja, zamke i modelski odgovori.

\section{Jedno upozorenje o izvornom tekstu zadatka}

Uz zadatak u repozitoriju postoji i prijepis
(\texttt{TASK/zadatak\_omrius.md}) u kojem uz svaki podzadatak stoji i numerirani popis
\emph{Koraka rješavanja}. Ti koraci \textbf{nisu} profesorov tekst --- dopisani su
naknadno kao pomoć pri radu. U ovom udžbeniku oni nikad ne ulaze u sivi okvir. Kad se
koriste, izričito su označeni kao izvedena uputa, a ne kao dio zadatka.
